\chapter{Soft Skill: Kemampuan Berkomunikasi}

\section{Identitas Mata Kuliah}
\begingroup
\renewcommand{\arraystretch}{1.15}
\noindent\begin{tabular}{@{}p{3.6cm}@{\hspace{0.2cm}}p{0.2cm}@{\hspace{0.2cm}}p{7.6cm}@{}}
    \textbf{Kode Mata Kuliah} & : & MII21-4010 \\
    \textbf{Nama Mata Kuliah} & : & Soft Skill: Kemampuan Berkomunikasi \\
    \textbf{Bobot} & : & 2 SKS \\
    \textbf{Bentuk Kegiatan} & : & Magang MBKM sebagai Software Engineer Intern di OPPO Manufacturing Indonesia \\
    \textbf{Proyek Terkait} & : & 360 Peer-to-Peer Grading System dan PQPR Software Project \\
\end{tabular}
\par\vspace{0.5\baselineskip}
\endgroup

Mata kuliah ini menekankan kemampuan berkomunikasi dalam lingkungan kerja profesional. Pada kegiatan magang, komunikasi menjadi bagian penting dalam memahami kebutuhan proyek, menyampaikan perkembangan pekerjaan, menjelaskan kendala teknis, menerima umpan balik, dan berkoordinasi dengan anggota tim yang memiliki peran berbeda.

\section{Silabus}
Silabus mata kuliah mencakup komunikasi lisan dan tertulis, penyampaian gagasan secara terstruktur, komunikasi dalam tim, kemampuan mendengarkan, pemberian dan penerimaan umpan balik, serta penyesuaian gaya komunikasi dengan lawan bicara. Dalam kegiatan magang, capaian tersebut diterapkan dalam interaksi dengan pembimbing lapangan, anggota tim teknis, pengguna internal, serta pihak yang memahami proses bisnis perusahaan.

Kemampuan komunikasi diperlukan karena pengembangan perangkat lunak tidak hanya dilakukan melalui aktivitas teknis. Setiap keputusan pengembangan perlu dijelaskan, setiap kebutuhan perlu dikonfirmasi, dan setiap perubahan perlu dikomunikasikan agar pekerjaan tim tetap selaras.

\section{Aktivitas}
Pada proyek 360 Peer-to-Peer Grading System, komunikasi dilakukan untuk memahami kebutuhan proses evaluasi karyawan dan memastikan alur sistem sesuai dengan tujuan pengguna. Diskusi dengan pembimbing dan anggota tim membantu memperjelas bentuk data yang perlu dikumpulkan, cara pengguna berinteraksi dengan sistem, serta informasi apa saja yang perlu ditampilkan dalam rekapitulasi penilaian.

Pada PQPR Software Project, komunikasi menjadi lebih penting karena proyek berkaitan dengan proses manufaktur yang memiliki istilah, aturan, dan kebutuhan operasional khusus. Penulis perlu memahami penjelasan dari pihak yang lebih mengetahui alur produksi, kemudian menerjemahkannya menjadi kebutuhan teknis yang dapat diterapkan dalam modul perangkat lunak. Proses ini membutuhkan kemampuan bertanya secara tepat, mencatat informasi penting, dan mengonfirmasi kembali pemahaman agar tidak terjadi kesalahan interpretasi.

Aktivitas komunikasi juga dilakukan melalui penyampaian progres pekerjaan, pelaporan kendala, diskusi perbaikan, dan penerimaan masukan terhadap fitur yang sedang dikembangkan. Dalam beberapa kondisi, kendala teknis perlu dijelaskan dengan bahasa yang mudah dipahami agar anggota tim nonteknis tetap dapat memahami dampaknya terhadap proses bisnis.

\section{Konsep Dasar yang Berkaitan dengan Silabus}
Konsep dasar yang digunakan adalah komunikasi dua arah. Komunikasi yang efektif tidak hanya berarti mampu menyampaikan pendapat, tetapi juga mampu mendengarkan, memahami konteks, dan memastikan pesan yang diterima sesuai dengan maksud pengirim. Konsep ini penting dalam kegiatan magang karena kebutuhan proyek sering kali berkembang setelah dilakukan diskusi dan evaluasi.

Konsep berikutnya adalah komunikasi lintas peran. Dalam proyek perangkat lunak, pengembang perlu berinteraksi dengan pembimbing, anggota tim teknis, pengguna, dan pihak operasional. Setiap pihak memiliki sudut pandang yang berbeda. Oleh karena itu, penyampaian informasi perlu disesuaikan dengan kebutuhan lawan bicara, misalnya menggunakan istilah teknis saat berdiskusi dengan tim pengembang dan menggunakan penjelasan proses saat berdiskusi dengan pengguna.

Kemampuan memberikan dan menerima umpan balik juga menjadi konsep penting. Umpan balik dari pembimbing dan pengguna membantu memperbaiki rancangan fitur, sedangkan kemampuan menyampaikan temuan atau usulan secara sopan dan terstruktur membantu proses pengambilan keputusan. Dengan komunikasi yang baik, perbaikan dapat dilakukan tanpa menghambat kerja sama tim.

\section{Hasil dan Pembahasan}
Hasil dari pelaksanaan mata kuliah ini adalah meningkatnya kemampuan berkomunikasi dalam konteks kerja profesional, terutama pada proyek perangkat lunak yang melibatkan kebutuhan teknis dan bisnis. Penulis belajar menyampaikan progres pekerjaan secara ringkas, menjelaskan kendala berdasarkan dampaknya terhadap pekerjaan, dan meminta klarifikasi ketika terdapat kebutuhan yang belum jelas.

Pada 360 Peer-to-Peer Grading System, komunikasi membantu memastikan bahwa fitur yang dikembangkan sesuai dengan kebutuhan proses evaluasi. Pada PQPR Software Project, komunikasi membantu menjembatani pemahaman antara kebutuhan operasional manufaktur dan implementasi teknis dalam sistem. Kedua pengalaman tersebut menunjukkan bahwa komunikasi berperan langsung terhadap ketepatan hasil pengembangan.

Pembahasan kegiatan menunjukkan bahwa kemampuan komunikasi menjadi faktor pendukung produktivitas tim. Kesalahan pemahaman dapat menyebabkan fitur tidak sesuai kebutuhan, sedangkan komunikasi yang jelas dapat mempercepat penyelesaian masalah. Dalam lingkungan magang, kemampuan untuk aktif bertanya, mencatat arahan, dan menindaklanjuti umpan balik menjadi kebiasaan kerja yang penting.

\section{Kesimpulan}
Mata kuliah Soft Skill: Kemampuan Berkomunikasi terlaksana melalui interaksi sehari-hari selama pengerjaan 360 Peer-to-Peer Grading System dan PQPR Software Project. Kegiatan ini memperkuat kemampuan menyampaikan ide, memahami kebutuhan pengguna, berdiskusi lintas peran, dan menerima umpan balik secara profesional. Kemampuan komunikasi terbukti menjadi bagian penting dari keberhasilan pengembangan perangkat lunak karena membantu memastikan pekerjaan teknis tetap sesuai dengan kebutuhan bisnis dan tujuan tim.
